\documentclass[conference]{IEEEtran}
\usepackage[T1]{fontenc}
\usepackage[utf8]{inputenc}
\usepackage{booktabs}
\usepackage{url}
\usepackage{array}
\usepackage{amsmath,amssymb}
\usepackage{microtype}
\usepackage{enumitem}
\usepackage[hidelinks]{hyperref}

\setlist[enumerate]{noitemsep,topsep=2pt,leftmargin=*}
\setlist[itemize]{noitemsep,topsep=2pt,leftmargin=*}

\title{Eleição em Anel e Two-Phase Commit para\\Transferências Bancárias Distribuídas}

\author{
  \IEEEauthorblockN{Integrante 1: \textit{[Nome Completo]}}
  \IEEEauthorblockA{RA: \textit{[Número de RA]}\\
  Instituto de Computação -- Unicamp\\
  MC714 -- Sistemas Distribuídos (2ºs/2026)}
  \and
  \IEEEauthorblockN{Integrante 2: \textit{[Nome Completo]}}
  \IEEEauthorblockA{RA: \textit{[Número de RA]}\\
  Instituto de Computação -- Unicamp\\
  MC714 -- Sistemas Distribuídos (2ºs/2026)}
  \and
  \IEEEauthorblockN{Entregáveis do Projeto}
  \IEEEauthorblockA{
    GitLab: \url{https://gitlab.ic.unicamp.br/mc714/2026s2/raXXXXXX}\\
    Vídeo: \url{https://youtu.be/exemplo-video-mc714}\\
    Demonstração em vídeo: $\approx$ 5 minutos
  }
}

\begin{document}
\maketitle

\begin{abstract}
Este trabalho apresenta a implementação e a avaliação experimental de dois algoritmos distribuídos clássicos e complementares: a eleição de líder em anel e o protocolo de compromisso atômico \textit{Two-Phase Commit} (2PC), integrados em uma aplicação de transferências bancárias replicadas. O sistema executa sobre contêineres OCI independentes orquestrados por Docker Compose v2, configurados em topologias de 3, 5 e 10 nós lógicos, comunicando-se exclusivamente por gRPC sobre TCP com prazos de expiração (\textit{deadlines}) explícitos. O líder eleito pelo anel assume o papel de coordenador do 2PC, enquanto os demais nós operam como coortes participantes com controle estrito de concorrência e travas (\textit{locks}) de contas em memória. A suíte automatizada de testes validou propriedades de segurança e vivacidade em fluxos normais de commit, abort por saldo insuficiente, tolerância a crash de participantes na Fase~1 e a ocorrência do clássico problema do bloqueio (\textit{blocking problem}) sob queda do coordenador pós-\textit{prepare}, demonstrando sua resolução operacional autônoma através da reeleição no anel. Os ensaios quantitativos evidenciaram latências médias de 2PC de 6,73~ms, 10,46~ms e 20,01~ms para 3, 5 e 10 nós, respectivamente, com detalhamento de percentis (p50, p95 e p99) e discussão da complexidade de mensagens.
\end{abstract}

\begin{IEEEkeywords}
sistemas distribuídos, eleição em anel, two-phase commit, gRPC, detecção de falhas, tolerância a falhas.
\end{IEEEkeywords}

\section{Introdução e Motivação}
A coordenação de processos autônomos e a garantia de atomicidade em operações de escrita constituem dois dos desafios mais fundamentais no projeto de sistemas distribuídos modernos~\cite{vsteen}. Em ambientes replicados onde múltiplas entidades compartilham a responsabilidade de manter a coerência de um estado global, a ausência de um coordenador central fixo exige protocolos robustos para a eleição dinâmica de líderes. Paralelamente, quando uma transação envolve mutações indivisíveis que abrangem múltiplos nós, protocolos de compromisso atômico tornam-se indispensáveis para assegurar que todos os nós convirjam para o mesmo desfecho (todos aplicam a alteração ou todos abortam), eliminando estados parciais inconsistentes.

Neste projeto, implementamos e investigamos empiricamente o Algoritmo de Eleição de Líder em Anel e o protocolo \textit{Two-Phase Commit} (2PC), integrando-os de maneira deliberada em uma aplicação prática de transferências financeiras entre réplicas contábeis. A escolha desse domínio bancário reduz a complexidade acessória de regras de negócio sem descaracterizar a propriedade central de interesse: a transferência entre contas (\texttt{conta\_A} e \texttt{conta\_B}) deve preservar estritamente a conservação do valor monetário total e a consistência mútua entre todas as réplicas ativas, garantindo isolamento através de travas (\textit{locks}) de conta.

A integração entre ambos os algoritmos resolve uma fragilidade inerente ao 2PC puro. No 2PC tradicional, o coordenador representa um ponto único de falha (\textit{single point of failure}); caso o coordenador colapse após os participantes votarem SIM na Fase~1 (\textit{prepare}) e antes de emitir a decisão da Fase~2, os participantes entram em estado de incerteza e bloqueio permanente, incapazes de decidir unilateralmente sem violar a segurança. Ao acoplar um detector de falhas baseado em \textit{heartbeats} periódicos e o algoritmo de eleição em anel, nosso sistema detecta a inatividade do coordenador, elege dinamicamente um novo líder e executa uma rotina cooperativa de resolução da pendência, restaurando a vivacidade do cluster.

O artefato foi desenvolvido em Python 3.10 com comunicação estritamente baseada em stubs gerados por \texttt{grpcio} v1.59.0 a partir de contratos Protocol Buffers (\texttt{proto3}), sendo orquestrado por contêineres Docker independentes. Conforme as diretrizes da disciplina, nenhum componente de terceiros que implemente os algoritmos distribuídos foi utilizado; toda a lógica do anel, coordenação do 2PC, transações contábeis e mecanismos de injeção de caos foi concebida e implementada integralmente pelos autores.

\section{Modelo do Sistema e Propriedades de Correção}
A especificação formal do modelo sob o qual os algoritmos operam é pré-requisito mandatório para delimitar garantias teóricas e interpretar resultados empíricos. A seguir, declaramos e justificamos os pressupostos para cada um dos dois algoritmos desenvolvidos.

\subsection{Algoritmo de Eleição de Líder em Anel}
\textbf{1. Modelo Temporal:} O modelo adotado é \textit{parcialmente síncrono}. Não se assume a existência de cotas superiores rígidas a priori para o tempo de transmissão em rede local; contudo, após estabilização, os prazos de expiração (\textit{deadlines} de RPC de 2,0~s) e os intervalos de \textit{heartbeat} (3,0~s com timeout de 1,5~s) fornecem limites práticos confiáveis para suspeitar da inatividade de vizinhos ou do líder, evitando bloqueios indefinidos.

\textbf{2. Propriedades dos Canais:} A comunicação utiliza gRPC sobre conexões TCP ponto a ponto, garantindo transporte ordenado (FIFO) e detecção de perdas por canal. Não há ordenação causal global entre conexões distintas. Chamadas expiradas retornam \texttt{DEADLINE\_EXCEEDED} ou \texttt{UNAVAILABLE}, sendo interpretadas como suspeita de inoperância.

\textbf{3. Tipos de Falha Considerados:} Considera-se falhas \textit{crash-stop} e \textit{crash-recovery}. O anel tolera a queda de nós intermediários através de salto dinâmico (\textit{skip}): ao detectar um sucessor imediato inativo, avança-se na topologia circular para o próximo nó disponível. Quedas do líder são diagnosticadas após 2 falhas consecutivas de heartbeat. Não são admitidas falhas bizantinas nem partições de rede desconexas.

\textbf{4. Hipóteses sobre Relógios, Persistência e Reinicialização:} Não se presume sincronização de relógios físicos; utilizam-se apenas timers monotônicos locais para timeouts. Identificadores pseudoaleatórios não sequenciais são gerados no boot. O anel opera em memória volátil, recalculando a topologia ativa a cada ciclo eleitoral.

\textbf{5. Propriedades de Segurança e Vivacidade:}
\begin{itemize}
  \item \textit{Segurança (Unicidade e Concordância):} Ao término de uma eleição, todos os nós alcançáveis concordam sobre o mesmo líder, rigorosamente o de maior identificador ativo:
  $\forall p, q \in \text{Vivos}, \; \text{Líder}_p = \text{Líder}_q = \max(\{ID_i\})$.
  \item \textit{Vivacidade (Término):} Qualquer eleição deflagrada em topologia com ao menos um nó ativo e caminhos alcançáveis eventualmente encerra com o anúncio do coordenador a todos os nós.
\end{itemize}

\subsection{Protocolo Two-Phase Commit (2PC) Bancário}
\textbf{1. Modelo Temporal:} Parcialmente síncrono. O coordenador adota deadlines gRPC de 2,5~s na Fase~1 (\textit{prepare}) e Fase~2. Qualquer expiração na Fase~1 é tratada conservadoramente como voto \textit{NÃO}, priorizando a segurança sobre a disponibilidade.

\textbf{2. Propriedades dos Canais:} RPCs ponto a ponto sobre gRPC/TCP. Os participantes mantêm histórico de transações decididas (\texttt{transaction\_history}), garantindo \textit{idempotência} frente a retentativas de rede.

\textbf{3. Tipos de Falha Considerados:} Tolera-se crash-stop de coortes na Fase~1 (resultando em abort seguro). Queda do coordenador pós-\textit{prepare} (o clássico \textit{blocking problem}) retém locks em \texttt{PREPARED} até resolução pelo novo líder. Nós recuperados sincronizam saldos via \texttt{SyncBankState}.

\textbf{4. Hipóteses sobre Relógios, Persistência e Reinicialização:} Sem relógios sincronizados. Travas locais (\texttt{bank\_lock}) gerenciam concorrência. Saldos confirmados e histórico são persistidos em snapshot JSON (\texttt{/tmp/bank\_state\_*.json}) na Fase~2, permitindo recuperação de estado após reinicialização.

\textbf{5. Propriedades de Segurança e Vivacidade:}
\begin{itemize}
  \item \textit{Segurança (Atomicidade e Consistência):} (i) Se uma réplica comita, todas as operacionais comitam; caso contrário, todas abortam; (ii) Nenhuma transação comita sem voto SIM unânime; (iii) A soma dos saldos de todas as contas é invariante ($\sum_{c} \text{Saldo}(c) = \text{constante}$); (iv) Locks exclusivos impedem disputas concorrentes sobre contas em \texttt{PREPARED}.
  \item \textit{Vivacidade (Progresso):} Transações sem falha de coordenador concluem em \texttt{COMMITTED} ou \texttt{ABORTED}. Sob crash pós-\textit{prepare}, o bloqueio é destravado assim que o anel elege um novo líder para conduzir a resolução.
\end{itemize}

\section{Arquitetura e Algoritmos}
Cada nó executa em contêiner independente (\texttt{src/node.py}) instanciando um servidor gRPC multithread (\texttt{ThreadPoolExecutor}) com os serviços \texttt{RingService} e \texttt{TwoPhaseCommitService}.

\subsection{Topologia e Descoberta de Vizinhos}
A ordem nominal do anel é passada por \texttt{RING\_TOPOLOGY} (ex: \texttt{no\_1:50051, \dots}). O nó determina sua lista circular de sucessores a partir de \texttt{MY\_ADDR}. Cada contêiner sorteia um identificador pseudoaleatório de 32 bits no boot, tornando a eleição não-trivial e desacoplada dos nomes dos nós.

\subsection{Algoritmo do Anel com Salto Dinâmico}
A eleição opera segundo a variante de van Steen e Tanenbaum~\cite{vsteen} com tolerância a falhas:
\begin{enumerate}
  \item \textbf{Deflagração:} O iniciador cria \texttt{ElectionMessage} com seu identificador e a transmite via \texttt{PassElection} ao primeiro sucessor responsivo.
  \item \textbf{Salto Dinâmico (\textit{Skip}):} Se o sucessor imediato não responder dentro do timeout (2,0~s), o nó remetente itera sobre os sucessores seguintes, saltando contêineres inativos até alcançar um nó vivo.
  \item \textbf{Fechamento de Ciclo:} Cada intermediário adiciona seu par \texttt{(id, addr)} e repassa a mensagem. Ao retornar ao iniciador, este elege o maior identificador do circuito.
  \item \textbf{Anúncio:} O iniciador difunde \texttt{AnnounceCoordinator} pelo anel. Cada participante atualiza seu líder e repassa a mensagem.
\end{enumerate}

\subsection{Detecção de Falhas via Heartbeat Periódico}
Em segundo plano, uma thread em cada nó seguidor emite a RPC \texttt{Ping} ao líder a cada 3,0~s (timeout 1,5~s). Se o líder falhar 2 vezes consecutivas (\texttt{MISSED\_HEARTBEATS\_LIMIT}), o nó assume a queda do líder e dispara nova eleição no anel.

\subsection{Protocolo Two-Phase Commit Bancário}
O líder atua como Coordenador do 2PC e os demais nós como coortes:
\begin{itemize}
  \item \textbf{Submissão:} O cliente invoca \texttt{SubmitTransaction} no líder (nós seguidores efetuam proxy transparente ao líder).
  \item \textbf{Fase 1 (Prepare):} O coordenador despacha paralelamente \texttt{PrepareTransaction}. As coortes validam contas, valor positivo, ausência de locks ativos e saldo suficiente. Em caso de aprovação, retêm os locks das contas em \texttt{pending\_transactions} e votam SIM; havendo falha ou timeout, votam NÃO.
  \item \textbf{Fase 2 (Commit/Abort):} Com votos SIM unânimes, o coordenador difunde \texttt{CommitTransaction}, efetivando débitos/créditos, liberando locks e gravando o snapshot local. Se houver voto NÃO ou timeout, difunde \texttt{AbortTransaction}, descartando o preparo e liberando as contas sem alterar saldos.
\end{itemize}

\subsection{Resolução Autônoma do Bloqueio Pós-Crash do Líder}
A injeção de crash no líder entre as fases retém coortes no estado \texttt{PREPARED}. O detector de heartbeat constata a queda e o anel elege um novo líder. O novo coordenador invoca \texttt{ResolveBlockedTransaction} com \texttt{ABORT}, liberando os locks de forma segura. Ao ser reiniciado, o antigo líder executa \texttt{SyncBankState} para sincronizar saldos com as réplicas ativas e purgar locks órfãos.

\section{Validação de Correção e Testes}
A correção foi comprovada por testes automatizados reproduzíveis cobrindo cenários normais e de falha para ambos os algoritmos.

\subsection{Correção da Eleição em Anel}
\begin{itemize}
  \item \textbf{Cenário Normal (Eleição em Topologia Íntegra):} Com 10 nós ativos, uma eleição foi disparada a partir do \texttt{no\_1}. A mensagem completou 10 saltos e o nó com maior identificador foi eleito e anunciado a 100\% dos nós, conforme verificado via \texttt{GetStatus}.
  \item \textbf{Cenário de Falha de Nó Intermediário (\textit{Skip}):} O contêiner \texttt{no\_3} foi parado (\texttt{docker stop no\_3}). Ao deflagrar eleição via \texttt{no\_2}, o nó detectou a indisponibilidade no timeout e saltou para \texttt{no\_4}, elegendo o líder correto entre os 9 nós vivos.
  \item \textbf{Cenário de Falha do Coordenador via Heartbeat:} O líder ativo foi encerrado. Após 6,0~s (2 heartbeats perdidos), os seguidores diagnosticaram a queda e elegeram autonomamente um novo líder.
\end{itemize}

\subsection{Correção do Two-Phase Commit Bancário}
A suíte \texttt{src/test\_2pc\_scenarios.py} foi executada sobre 10 nós em 11/11/2026, com 100\% de êxito em 18,60~s:
\begin{enumerate}
  \item \textbf{Commit Normal:} Transferência válida de R\$100,00 de \texttt{conta\_A} para \texttt{conta\_B}; todas as 10 réplicas confirmaram saldos de R\$800,00 e R\$1200,00, sem travas residuais e com soma total invariante.
  \item \textbf{Abort por Regra de Negócio:} Tentativa de transferência de R\$51.000,00 (saldo insuficiente); coortes emitiram voto NÃO na Fase~1, resultando em abort global com saldos e locks preservados.
  \item \textbf{Crash de Participante na Fase 1:} Coorte \texttt{no\_4} pausada via \texttt{SetChaosHook}; timeout no coordenador tratado como voto NÃO, cancelando a operação sem alterar saldos dos sobreviventes.
  \item \textbf{Crash do Coordenador pós-Prepare (Bloqueio):} Queda induzida no líder pós-coleta de votos SIM; coortes comprovaram o bloqueio teórico em \texttt{PREPARED}. O anel elegeu novo líder, que acionou \texttt{ResolveBlockedTransaction}, liberando os locks com sucesso. O nó caído reiniciou e sincronizou saldos via \texttt{SyncBankState}.
\end{enumerate}

\section{Metodologia Experimental e Ambiente}
Avaliou-se a escalabilidade dos dois algoritmos manipulando o tamanho do cluster ($N$) em três níveis: \textbf{3, 5 e 10 nós lógicos} (\texttt{docker/compose-3nodes.yaml}, \texttt{compose-5nodes.yaml} e \texttt{compose.yaml}).

As métricas analisadas foram: (i) \textbf{Tempo de Convergência da Eleição (ms)}, do disparo da eleição até o encerramento do anúncio; e (ii) \textbf{Latência do 2PC (ms)}, do envio do \textit{prepare} até a confirmação da Fase~2. Executaram-se 20 rodadas independentes ($R=20$) por topologia, com alternância de contas e valores (R\$1 a R\$10) e intervalo de 0,3~s entre transações.

\textbf{Ambiente de Execução:} Processador AMD Ryzen 5 5600 6-Core (4 vCPUs alocadas, 3,8~GiB RAM DDR4), Debian Linux (Kernel 6.12.116), Docker Engine 29.9.3, Compose v2, imagem \texttt{python:3.10-slim}, \texttt{grpcio 1.59.0}. Duração agregada dos ensaios de $\approx 110$~s sem aquecimento prévio separado. Estatísticas (Média, DP, p50, p95 e p99) obtidas por interpolação linear.

\section{Resultados e Discussão}
Os dados experimentais consolidados constam nas Tabelas~\ref{tab:ring} e~\ref{tab:2pc}.

\begin{table}[t]
\caption{Convergência da eleição em anel (ms, 20 rodadas).}
\label{tab:ring}
\centering
\small
\begin{tabular}{crrrrrr}
\toprule
Nós & Sucesso & Média & DP & p50 & p95 & p99 \\
\midrule
3  & 20/20 & 107,57 & 2,05 & 107,00 & 110,77 & 111,45 \\
5  & 20/20 & 107,65 & 1,35 & 107,31 & 109,08 & 109,80 \\
10 & 20/20 & 107,56 & 1,23 & 107,23 & 108,50 & 109,20 \\
\bottomrule
\end{tabular}
\end{table}

\begin{table}[t]
\caption{Latência ponta a ponta do 2PC (ms, 20 commits).}
\label{tab:2pc}
\centering
\small
\begin{tabular}{crrrrrr}
\toprule
Nós & Sucesso & Média & DP & p50 & p95 & p99 \\
\midrule
3  & 20/20 & 6,73 & 0,78 & 6,52 & 7,68 & 8,99 \\
5  & 20/20 & 10,46 & 1,21 & 10,17 & 12,83 & 14,02 \\
10 & 20/20 & 20,01 & 2,16 & 19,31 & 24,18 & 26,16 \\
\bottomrule
\end{tabular}
\end{table}

\subsection{Análise do Protocolo 2PC Bancário}
A Tabela~\ref{tab:2pc} evidencia um aumento estrito da latência média com o número de coortes: 6,73~ms (3 nós), 10,46~ms (5 nós) e 20,01~ms (10 nós), um salto de quase $3\times$ entre 3 e 10 nós. Para $N$ nós, o 2PC requer $2(N-1)$ mensagens entre coordenador e coortes ($O(N)$ mensagens). Como a Fase~1 e a Fase~2 aguardam a última resposta no caminho crítico, o acréscimo de participantes eleva a probabilidade de atrasos por chaveamento de threads e contenção de CPU, expandindo a cauda de latência: o p99 aumentou de 8,99~ms para 26,16~ms. A taxa de êxito foi de 100\% nas 60 transações.

\subsection{Análise da Eleição em Anel e Limitações de Medição}
Na Tabela~\ref{tab:ring}, as médias permaneceram estabilizadas em $\approx 107,6$~ms em todas as escalas. Embora a teoria dite $2N$ encaminhamentos ($O(N)$ saltos), o script cliente consulta o status por amostragem periódica de 100~ms (\texttt{time.sleep(0.1)}). Como a propagação das mensagens gRPC sobre a rede loopback local consome menos de 10~ms mesmo em 10 nós, o valor registrado reflete a granularidade do \textit{polling} somada ao handshake gRPC, mascarando a complexidade assintótica do algoritmo e demonstrando a relevância de registrar carimbos temporais internos aos nós.

\section{Ameaças à Validade e Limitações}
\begin{itemize}
  \item \textbf{Contêineres em Host Único:} Os nós compartilham a mesma CPU física, memória e kernel Linux, isolando processos mas não constituindo domínios de falha independentes.
  \item \textbf{Persistência e Semântica de Recuperação:} O snapshot de saldos em JSON permite recuperação rápida de saldos, mas transações em trânsito e locks residem em memória volátil, sem Write-Ahead Log com \texttt{fsync} síncrono.
  \item \textbf{Partições Assimétricas:} Quedas prolongadas que fragmentem o anel paralisam a circulação de mensagens na ausência de mecanismos de quórum.
\end{itemize}

\section{Itens Extras e Atividades de Bônus}
Conforme a Seção~5 do regulamento, o projeto implementou três extensões avançadas:
\textbf{1) Integração Coerente em Aplicação Prática:} O líder do anel coordena o 2PC e a reeleição é usada ativamente para destravar transações bloqueadas pós-queda;
\textbf{2) Terceiro Protocolo (Detecção de Falhas por Heartbeat):} Detector autônomo baseado em RPCs \texttt{Ping} periódicas com contadores de perda configuráveis;
\textbf{3) Mecanismo de Crash-Recovery e Reconciliação:} Implementação de \texttt{SyncBankState} para restauração de saldos e eliminação de travas órfãs em nós reiniciados.

\section{Declaração de Autoria e Código Reutilizado}
Toda a arquitetura, lógica dos algoritmos de anel com salto dinâmico (\textit{skip}), protocolo 2PC, transações contábeis, heartbeat, testes e scripts experimentais foram concebidos e implementados integralmente pelos membros da equipe. Utilizou-se unicamente bibliotecas padrão do Python 3.10 e \texttt{grpcio}/\texttt{grpcio-tools} v1.59.0 para transporte RPC e geração de stubs dos arquivos \texttt{.proto}. Nenhuma biblioteca de consenso ou commit atômico de terceiros foi empregada.

\section{Instruções de Reprodução}
Todo o ambiente é reprodutível via Docker Compose v2: inicialização com \texttt{docker compose up -d --build}; validação com \texttt{docker compose exec no\_1 python src/test\_2pc\_scenarios.py}; coleta experimental com \texttt{src/run\_experiments.py} e \texttt{src/run\_2pc\_experiments.py} (\texttt{--rounds 20}); topologias com \texttt{-f docker/compose-3nodes.yaml} ou \texttt{5nodes.yaml}; e encerramento com \texttt{docker compose down -v}.

\section{Conclusão}
A integração entre eleição em anel e 2PC permitiu materializar o compromisso clássico entre disponibilidade e atomicidade estrita em sistemas distribuídos. O detector de falhas e a reeleição proporcionam uma recuperação operacional prática para o problema do bloqueio, enquanto as medições quantitativas evidenciaram o impacto da escala sobre a cauda de latência do 2PC e a necessidade de instrumentação de alta resolução em avaliações de sistemas concorrentes.

\begin{thebibliography}{00}
\bibitem{vsteen} M. van Steen and A. S. Tanenbaum, \textit{Distributed Systems}, 4th ed., distributed-systems.net, 2023.
\bibitem{grpc} gRPC Authors, ``gRPC Python Documentation,'' 2026. [Online]. Disponível em: \url{https://grpc.io/docs/languages/python/}. Acesso em: 11-Nov-2026.
\bibitem{docker} Docker Inc., ``Docker Compose Documentation v2,'' 2026. [Online]. Disponível em: \url{https://docs.docker.com/compose/}. Acesso em: 11-Nov-2026.
\end{thebibliography}

\end{document}
